\section{Discussion}
\label{sec:discussion}

KM shows that completion, consequence-based saturation, and hypertableau search
can share one frontend, one clause representation, and one contract for
publishing an answer, and that a router over source features then selects among
them without access to the expected result.  On held-out inputs of the
ORE~2015 pool this combination solves as many inputs as Konclude within the
same limits, in less time and memory on the inputs both solve.  The
consequence-based procedures answer most inputs, and the hypertableau
procedures answer the inputs on which saturation exhausts its budget, which
include GALEN and large $\mathcal{SRIQ}$ terminologies.  \todo{Add the
per-route counts of the v1.7.0 pool run.}

Three design decisions distinguish KM from the reasoners it builds on.  First,
every procedure returns a complete answer or declines.  A fast procedure for a
fragment can therefore be added without a proof that the router sends it only
inputs of that fragment, because the procedure itself checks the clauses it
receives.  Second, a procedure runs in its own process, which bounds the memory
of a classification by the largest worker and makes a time or memory limit
enforceable from outside.  Third, the evidence for an answer can be checked
by a program derived from a Lean proof.  The checked model for a negative
answer is the component that scales: it certifies 1,618 of 1,696 pool answers
without a completeness proof for any optimised procedure.

The port of Konclude's kernel deserves a separate remark.  Konclude's
saturation, caching, and classification strategies are the result of a decade
of engineering, and reimplementing them from the published descriptions would
have produced a different and weaker system.  Translating the code function by
function preserved that engineering, and the differential tests then found two
defects that the port had inherited and that the original also shows.  The
port is distributed under Konclude's licence, the LGPL~v3.

\subsection{Limitations}
\label{sec:limitations}

The evaluation corpus is ORE~2015.  The held-out half is disjoint from the
inputs used for development decisions, and it is drawn from the same
collection as the design half, so it measures generalisation within that
collection.  A snapshot of 189 current OBO Foundry ontologies was classified
with KM~v1.3.0 only; that release completed 97 of them against 139 for
Konclude, mostly because its frontend refused inverse properties in positions
that later releases accept.  Repeating this measurement with v1.7.0, and
adding SNOMED~CT and the large OBO ontologies with imports merged, is the next
evaluation step.

Twelve admitted inputs of the pool remain unsolved within 240 seconds.  Inputs
in which many individuals may be equal to the same nominal can take minutes,
and some property-value queries on such inputs exceed the plugin's time limit.
KM reads functional syntax, OWL/XML, RDF/XML, and Turtle; it does not read
Manchester syntax, and it performs no network access, so imports must be
merged into one document before classification.

The Lean results cover the reasoning rules, the checkers, and reference
procedures on fragments.  The optimised Rust procedures, the parser, and the
router are validated by the checkers on individual answers and by tests; about
5\% of the pool answers have no certificate because a checker declined or
exceeded its budget.  Differential testing found wrong answers in every release
line before v1.7.0, and the absence of a wrong answer in the v1.7.0 campaigns
is evidence for the tested distributions of ontologies.

\subsection{Availability}
\label{sec:availability}

KM is developed at
\url{https://github.com/bio-ontology-research-group/kobayashi-marust}.  The
release \texttt{v1.7.0} provides the \texttt{km} executable for Linux and macOS
on x86-64 and ARM and for Windows on x86-64, and the plugin archives for
Prot\'eg\'e~5.6 and Prot\'eg\'e~6.  The Rust and Lean sources and the plugin are
released under the BSD 3-Clause licence, and the Konclude port under the
LGPL~v3.  The repository contains the benchmark harness, the per-input result
files of the measurements reported here, the errata with reproducers, and the
scripts that run the Lean gates.  \todo{Archive the v1.7.0 release and the
benchmark records with a DOI and cite it here.}

\section{Conclusion}

\KM{} combines the three families of OWL~2~DL classification procedures in one
Rust executable, selects among them from features of the source, and publishes
only complete answers.  It answers class, individual, and property queries,
explains entailments, updates a classification after an edit, and runs inside
Prot\'eg\'e.  On held-out ORE~2015 ontologies it solves as many inputs as
Konclude with lower time and memory on common inputs, and Lean checkers certify
95\% of its answers on the full pool.
